\subsection{截切序列的构造}

命题 9.--- 设 A 为诺特环，M 为有限生成 A-模，a 是关于加法和乘法稳定的 A 的子集，\(\mathfrak{q}_{1},\ldots,\mathfrak{q}_{m}\) 是不包含 a 的 A 的素理想，且 \(r\geqslant1\) 为满足下式的整数：

\[
\operatorname{codim} (\mathrm{V} (\mathfrak {a}) \cap \operatorname{Supp} (\mathrm{M}), \operatorname{Supp} (\mathrm{M})) \geqslant r.\tag{21}
\]

存在一个由 a 中元素组成的序列 \((x_{1}, \ldots, x_{r})\)，它不属于理想 \(q_{1}, \ldots, q_{m}\) 中的任何一个，并且序列 \((x_{1}/1, \ldots, x_{r}/1)\) 由 \(A_{p}\) 中元素组成，是 \(A_{p}\)-模 \(M_{p}\) 的截切序列，对每个满足 \(p \in V(a) \cap \text{Supp}(M)\) 的素理想都成立。

我们对 \(r\) 归纳论证。首先设 \(r = 1\)，并以 \(\Phi\) 表示 \(\operatorname{Supp}(M)\) 的（有限）极小元集合。反设不存在满足陈述条件的元素 \(x_1\)（属于 \(\mathfrak{a}\)）。取 \(x \in \mathfrak{a}\)，且不属于 \(\mathfrak{q}_1 \cup \ldots \cup \mathfrak{q}_m\)。根据假设，存在一个素理想 \(\mathfrak{p}\)，属于 \(V(\mathfrak{a}) \cap \operatorname{Supp}(M)\)，使 \(x/1\)（即 \(x\) 在 \(A_{\mathfrak{p}}\) 中的像）不是 \(M_{\mathfrak{p}}\) 的截切元。令 \(\Psi\) 为由素理想 \(\mathfrak{q} \in \Phi\) 且包含于 \(\mathfrak{p}\) 者组成的集合；则 \(\operatorname{Supp}(M_{\mathfrak{p}})\) 的极小元是素理想 \(\mathfrak{q}A_{\mathfrak{p}}\)，它们属于 \(A_{\mathfrak{p}}\)，其中 \(\mathfrak{q}\) 遍历 \(\Psi\)。由 no. 2 的 prop. 3，存在 \(\mathfrak{q} \in \Psi\)，使 \(x/1 \in \mathfrak{q}A_{\mathfrak{p}}\)，因而 \(x \in \mathfrak{q}\)。换言之，有 \(\mathfrak{a} \subset \mathfrak{q}_1 \cup \ldots \cup \mathfrak{q}_m \cup \bigcup_{\mathfrak{q} \in \Phi} \mathfrak{q}\)。由于有 \(\mathfrak{a} \not\subset \mathfrak{q}_j\) 对 \(1 \leqslant j \leqslant m\) 成立，II, § 1, no. 1 的 prop. 2 证明存在素理想 \(\mathfrak{q} \in \Phi\)，包含 \(\mathfrak{a}\)，因而

\[
\mathrm{V} (\mathfrak {q}) \subset \mathrm{V} (\mathfrak {a}) \cap \operatorname{Supp} (\mathrm{M}).
\]

由于 V(a) 包含 Supp(M) 的一个不可约分支，这与假设矛盾：

\[
\operatorname{codim} (\mathrm{V} (\mathfrak {a}) \cap \operatorname{Supp} (\mathrm{M}), \operatorname{Supp} (\mathrm{M})) \geqslant 1.
\]

现在设 \(r \geqslant 2\)。根据归纳假设，可以找到一个序列 \((x_1, ..., x_{r-1})\)，其元素属于 \(\mathfrak{a}\)，不属于 \(\mathfrak{q}_1 \cup ... \cup \mathfrak{q}_m\)，并且对于每个 \(\mathfrak{p} \in V(\mathfrak{a}) \cap \operatorname{Supp}(M)\)，序列 \((x_1/1, ..., x_{r-1}/1)\) 由 \(A_{\mathfrak{p}}\) 中元素组成，是 \(M_{\mathfrak{p}}\) 的截切序列。置 \(N = M / \sum_{i=1}^{r-1} x_i M\)。只需构造元素 \(x_r\)，属于 \(\mathfrak{a}\)，不属于 \(\mathfrak{q}_1 \cup ... \cup \mathfrak{q}_m\)，并且对于每个 \(\mathfrak{p} \in V(\mathfrak{a}) \cap \operatorname{Supp}(M)\)，其像 \(x_r/1\) 属于 \(A_{\mathfrak{p}}\)，且它是 \(N_{\mathfrak{p}}\) 的截切元。根据证明的第一部分，只需证明以下两个关系：

\begin{enumerate}
\def\labelenumi{(\arabic{enumi})}
\setcounter{enumi}{21}
\tightlist
\item
\end{enumerate}

\[
\mathrm{V} (\mathfrak {a}) \cap \operatorname{Supp} (\mathrm{M}) = \mathrm{V} (\mathfrak {a}) \cap \operatorname{Supp} (\mathrm{N}),\tag{23}
\]

\[
\operatorname{codim} (\mathrm{V} (\mathfrak {a}) \cap \operatorname{Supp} (\mathrm{N}), \operatorname{Supp} (\mathrm{N})) \geqslant 1.
\]

现在有

\[
\operatorname{Supp} (\mathrm{N}) = \mathrm{V} (x _ {1}) \cap \dots \cap \mathrm{V} (x _ {r - 1}) \cap \operatorname{Supp} (\mathrm{M})
\]

根据 II, § 4, no. 4 的 prop. 18 的推论，并且由于 \(x_1\)、\ldots、\(x_{r-1}\) 属于 a，可得 (22)。不等式 (23) 随后由假设 (21) 以及 no. 3 的 prop. 4, a) 得到，其中取

\[
m = r - 1, \quad X = \operatorname{Supp} (M), \quad Y = V (\mathfrak {a}) \cap \operatorname{Supp} (N), \quad H _ {i} = V (x _ {i})
\]

从而 X' = Supp(N)。

推论 1. --- 设 A 为诺特环，M 为有限生成 A-模，\(\mathfrak{p}_1, \ldots, \mathfrak{p}_n\)、\(\mathfrak{q}_1, \ldots, \mathfrak{q}_m\) 为 A 的素理想，且 \(r\) 为整数 \(\geqslant 1\)。假设有 \(\mathfrak{p}_i \notin \mathfrak{q}_j\)，其中 \(1 \leqslant i \leqslant n\) 且 \(1 \leqslant j \leqslant m\)，并且有 \(\dim_{\mathbf{A}_{\mathfrak{p}_i}} (\mathbf{M}_{\mathfrak{p}_i}) \geqslant r\) 对 \(1 \leqslant i \leqslant n\) 成立。则存在一个由 A 中元素组成的序列 \((x_1, \ldots, x_r)\)，它属于所有 \(\mathfrak{p}_i\) 而不属于任何 \(\mathfrak{q}_j\)，并且对于 \(1 \leqslant i \leqslant n\)，\(x_1, \ldots, x_r\) 在 \(\mathbf{A}_{\mathfrak{p}_i}\) 中的像组成 \(M_{\mathfrak{p}_i}\) 的截切序列。

置 \(\mathfrak{a} = \bigcap_{i}\mathfrak{p}_{i}\)。由于 \(\mathbf{M}_{\mathfrak{p}_i}\neq \{0\}\)，有 \(\mathfrak{p}_i\in \mathrm{Supp}(\mathbf{M})\)（\(1\leqslant i\leqslant n\)），从而 \(\mathrm{V}(\mathfrak{a})\subset \mathrm{Supp}(\mathbf{M})\)。我们有

\[
\begin{array}{r l} \operatorname{codim} (\mathrm{V} (\mathfrak {a}) \cap \operatorname{Supp} (\mathrm{M}), \operatorname{Supp} (\mathrm{M})) & = \operatorname{codim} (\mathrm{V} (\mathfrak {a}), \operatorname{Supp} (\mathrm{M})) = \\ & = \inf _ {i} \left(\operatorname{codim} (\mathrm{V} (\mathfrak {p} _ {i}), \operatorname{Supp} (\mathrm{M}))\right) = \inf _ {i} \dim (\mathrm{M} _ {\mathfrak {p} _ {i}}) \geqslant r \end{array}
\]

（§ 1, no. 4, prop. 9），于是应用 prop. 9。

推论 2. --- 设 A 为诺特局部环，M 为非零有限生成 A-模，a 是 \(\mathfrak{m}_{\mathbf{A}}\) 的一个关于加法和乘法稳定的子集，并满足 long(M/aM) \textless{} + ∞。a 中存在一个对于 M 为极大截切的子集。

事实上，有 \(\mathrm{V}(\mathfrak{a}) \cap \operatorname{Supp}(\mathbf{M}) = \operatorname{Supp}(\mathbf{M}/\mathfrak{a}\mathbf{M}) = \{\mathfrak{m}_{\mathbf{A}}\}\)（§ 1, no. 4, remark 1），因而 \(\operatorname{codim}(\mathrm{V}(\mathfrak{a}) \cap \operatorname{Supp}(\mathbf{M}), \operatorname{Supp}(\mathbf{M})) = \dim(\operatorname{Supp}(\mathbf{M})) = \dim_{\mathbf{A}}(\mathbf{M})\)，于是应用 prop. 9。

